\section{Introduction}\label{sec:intro}
Given digraphs $H$ and $G$, an \emph{$H$-factor} in $G$ is a collection of vertex-disjoint copies of $H$ in $G$ which together cover every vertex of $G$.
Note that $G$ can only have an $H$-factor if $|H|$ divides $|G|$, and a natural question, in the spirit of Dirac's theorem, is to determine which minimum degree condition guarantees an $H$-factor when this divisibility condition holds.
For digraphs, the natural degree parameter is the \emph{minimum semidegree} $\delta^0(G)$ of $G$, which is the minimum of its minimum outdegree $\delta^+(G)$ and its minimum indegree $\delta^-(G)$.
Throughout this paper, digraphs are finite, loopless and have no multiple edges in the same direction, but they may contain pairs of opposite edges, that is, $2$-cycles. The copies of $H$ in an $H$-factor need not be induced.

In this paper, we consider the case where $H$ is an orientation of a cycle, that is, a cycle each of whose edges is given one of the two directions.
Molla and Treglown~\cite[Theorem~1.7]{MT} proved that, for fixed $r$ and every $\eta>0$, every sufficiently large digraph $G$ on $n$ vertices with $r\mid n$ and $\delta^0(G)\ge(1/2+1/(2r)+\eta)n$ has an $H$-factor for every orientation $H$ of a cycle of length $r$.
For orientations with equally many forward and backward edges, their result on the total degree gives the coefficient $1/2$ in place of $1/2+1/(2r)$~\cite[Corollary~1.3]{MT}.
They asked whether the coefficient $1/2$ suffices for every orientation of an even cycle~\cite[Question~7.2]{MT}, and, in particular, for the directed four-cycle~\cite[Conjecture~7.1]{MT}.
We show that this is the case.

\begin{theorem}\label{compact390:main}
For every orientation $H$ of an even cycle of length $r\ge4$
and every $\eta>0$, there is some $n_0=n_0(H,\eta)$ such that every
digraph $G$ with
\[
 |G|=n\ge n_0,\qquad r\mid n,\qquad
 \delta^0(G)\ge(1/2+\eta)n
\]
has an $H$-factor.
\end{theorem}

The coefficient $1/2$ in Theorem~\ref{compact390:main} cannot be reduced. Indeed, take two vertex-disjoint complete digraphs whose orders are nearly equal and not divisible by $r$, but whose total order $n$ is divisible by $r$. Its minimum semidegree is $n/2-O(r)$. Since $H$ is connected, every copy of $H$ lies in one of the two complete digraphs, whose order is not divisible by $r$; hence there is no $H$-factor.

\smallskip

\noindent\textbf{Fractional tilings.}
As is common for factor problems, the main difficulty lies in a fractional version of the problem, which we now describe.
A \emph{homomorphism} from $H$ to a digraph $F$ is a map $f:V(H)\to V(F)$ such that $f(x)f(y)$ is an edge of $F$ whenever $xy$ is an edge of $H$. A homomorphism need not be injective, so it may identify nonadjacent vertices of $H$, and we write $m_f(v)=|f^{-1}(v)|$ for the number of vertices of $H$ which $f$ maps to $v$.
A \emph{perfect fractional homomorphic $H$-tiling} of $F$ assigns a weight $\lambda_f\geq 0$ to each homomorphism $f$ from $H$ to $F$, so that $\sum_f m_f(v)\lambda_f=1$ for every vertex $v$ of $F$.
We prove the following theorem.

\begin{theorem}\label{compact390:finite}
Let $H$ be an orientation of an even cycle of length at least four. Then, every finite digraph $F$ with $\delta^0(F)>|F|/2$ has a perfect
fractional homomorphic $H$-tiling.
\end{theorem}

Theorem~\ref{compact390:finite} implies Theorem~\ref{compact390:main} by standard applications of the regularity and absorption methods, and, in fact, the two statements are equivalent (see Theorem~\ref{general:equivalence}).
Note also that Theorem~\ref{compact390:finite} concerns digraphs $F$ of every order, and has no error term in its condition $\delta^0(F)>|F|/2$.

\subsection{Methods}\label{sec:methods}
By linear programming duality, Theorem~\ref{compact390:finite} says that every \emph{fractional $H$-cover} of $F$, that is, every nonnegative weighting of $V(F)$ under which each homomorphism from $H$ has total weight at least one, has total weight at least $|F|/|H|$.
We therefore fix a cover of smaller total weight and seek a contradiction.
The condition $\delta^0(F)>|F|/2$ enters through a simple observation. If we delete one vertex from $H$ and map the resulting path into $F$, then the vertices which can be put at the deleted position form the intersection of an outneighbourhood or inneighbourhood with another such neighbourhood. Both have more than $|F|/2$ vertices, so they intersect, and so the path extends to a homomorphism of $H$. The whole proof consists of weighted versions of this observation.

How the observation is used depends on the difference $k_H$ between the numbers of forward and backward edges of $H$.
When $k_H=0$ or $k_H\ge4$, and in many cases when $k_H=2$, we can map $H$ onto directed four-cycles and six-cycles so evenly that a fractional tiling of $F$ by directed four-cycles and six-cycles gives one by $H$. For directed cycles, and for a few short orientations, we delete a vertex of $H$ and look at the cheapest copies of the resulting path. We show that, for a cover of small weight, one of them has two completion neighbourhoods of small total weight under a suitable measure. This contradicts the semidegree condition by a double count.
The remaining orientations have $k_H=2$, and form two families. For these, we fix a vertex of weight zero, or an edge of minimum weight. We then use homomorphisms of $H$ into a few vertices around it to bound from below the weight of every vertex in the neighbourhoods of these vertices. As each of these neighbourhoods has more than $|F|/2$ vertices, these bounds force the total weight to be at least $|F|/|H|$.
We explain this in more detail in Section~\ref{sec:overview}.

\smallskip

\noindent\textbf{Organisation of the paper.}
Section~\ref{sec:overview} outlines the proof and deduces Theorem~\ref{compact390:finite} from the results for each family.
Section~\ref{framework} collects the tools used throughout: linear programming duality, the deletion of vertices of maximum weight, and bounds on weights in neighbourhoods.
Section~\ref{compact390:structural-section} reduces the problem to directed four-cycles and six-cycles and two families of orientations.
In Section~\ref{norm391:section}, we analyse cheapest paths whose completion neighbourhoods are heavy, and deduce bounds on optimal covers which we use for the two families.
We treat directed four-cycles and six-cycles in Section~\ref{sec:directed}, the first family in Sections~\ref{sec:spantwo} and~\ref{height390:three-section}, and the second family in Section~\ref{sec:tworeturns}.
Finally, Section~\ref{factor402:section} passes from fractional tilings to factors.
